\Needspace{18\baselineskip}
\section*{AI use}\label{sec:ai-use}
The research for this paper was carried out with autonomous rStation agents
using Gemini 4 Argon and GPT-6 Astra. The author launched these agents, which
conducted the research without human intervention. The initial arguments
were recorded as
\href{https://rstation.ai/problems/tenth_discriminant/papers/2/v12}{paper 2,
version 12} (under the agent name Severi Jet), giving the bound $499$, and
\href{https://rstation.ai/problems/tenth_discriminant/papers/22/v1}{paper 22,
version 1} (under the agent name Mira Cairn), giving $498$.
The local arithmetic, resultant and global estimates used here are in
Sections 3--6 of the former; the spectral estimates are developed in
\href{https://rstation.ai/problems/tenth_discriminant/papers/16/v2}{paper 16,
version 2}, Sections 2--5 and 6.4. The necessary derivations are included in the appendices. The proof was subsequently checked externally by separate
AI agents using GPT-6 Astra and Claude Opus 5.5, run independently of the
agents that made the discoveries.

GPT-6 Astra, Claude Opus 5.5 and Gemini 4 Argon were then used together to
simplify the proof and prepare this exposition. This included revisions of
the outline and individual sections, mathematical and editorial review,
and recomputation of the numerical bounds. In particular, narrowing the
auxiliary-prime interval gives the shorter rank argument proved in
Section~\ref{sec:rank}.
